\subsection{Higher-Order Tables and Problem Patterns}
\label{subsec:higher-order-tables-and-problem-patterns}
% Source integration: attached Chapter 13 screenshots, Sections 13.7--13.8.

Two-way tables are the simplest categorical-data structure.  Higher-order tables introduce several possible notions of dependence: full independence, conditional independence given a third factor, marginal homogeneity, and interactions among factors.

\begin{definition} \label{def:two-way-log-linear-independence} \hlt{Two-Way Log-Linear Form}\\
For a two-way table with cell probabilities $\pi_{ij}$, independence is
\[
\pi_{ij}=\pi_{i\cdot}\pi_{\cdot j}.
\]
Equivalently,
\[
\log\pi_{ij}
=
\log\pi_{i\cdot}+\log\pi_{\cdot j}.
\]
Writing $\alpha_i=\log\pi_{i\cdot}$ and $\beta_j=\log\pi_{\cdot j}$, departures from independence may be represented by interaction terms:
\[
\log\pi_{ij}=\alpha_i+\beta_j+\gamma_{ij}.
\]
\end{definition}

\begin{remark} \label{rem:three-way-log-linear-models} \hlt{Three-Way Interactions}\\
For a three-way table with factors indexed by $i,j,k$, one possible log-linear model is
\[
\log\pi_{ijk}
=
\alpha_i+\beta_j+\delta_{ij}+\varepsilon_{ik}+\gamma_{jk}.
\]
This form permits pairwise, second-order dependencies but excludes a third-order interaction.  Likelihood-ratio tests can compare nested log-linear models by counting the parameters removed, as in Theorem \ref{thm:wilks-likelihood-ratio-chi-square}.
\end{remark}

\begin{principle} \label{prin:categorical-data-design-before-test} \hlt{Choose the Test from the Sampling Mechanism}\\
Before applying a categorical-data test, identify what was fixed by design and what was random:
\begin{enumerate}[label=(\roman*)]
\setlength{\itemsep}{0pt}
\item fixed margins in a $2\times2$ randomized or conditional analysis suggest Fisher's exact test;
\item independent multinomial samples suggest a homogeneity test;
\item one cross-classified multinomial sample suggests an independence test;
\item matched pairs suggest McNemar's test or a matched-table extension;
\item case-control sampling can estimate odds ratios but not absolute disease risks without extra information.
\end{enumerate}
The same displayed table can therefore support different questions under different sampling stories.
\end{principle}

\begin{remark} \label{rem:counts-not-percentages-in-chi-square-tests} \hlt{Use Counts, Not Percentages}\\
Pearson's statistic is calibrated for counts.  If every observed and fitted count in a table is divided by $n$, the resulting statistic is divided by $n$:
\[
\sum_i
\frac{(O_i/n-E_i/n)^2}{E_i/n}
=
\frac1n\sum_i\frac{(O_i-E_i)^2}{E_i}.
\]
Using percentages instead multiplies this relative-frequency statistic by $100$.  Thus a chi-square statistic computed from relative frequencies or percentages is not the same statistic unless the original count scale is restored.
\end{remark}

\begin{qnbank} \label{qnbank:categorical-data-source-problem-patterns} \hlt{Source Problem Patterns}\\
The source problems attached to this chapter provide data sets and prompts for several recurring categorical-data tasks:
\begin{enumerate}[label=(\roman*)]
\setlength{\itemsep}{0pt}
\item testing homogeneity or independence in genetic, mortality, blood-group, classroom-grade, treatment, and attitude tables;
\item distinguishing independence from homogeneity when the same numerical table could arise from different sampling schemes;
\item extending two-way ideas to three factors, as in education, age, and political interest;
\item using Fisher's exact test for small $2\times2$ tables;
\item applying McNemar's test to matched exposure windows, before--after categorical outcomes, or paired controls;
\item estimating and interpreting odds ratios in case-control and prospective designs.
\end{enumerate}
The screenshots supply the problem statements and data tables, but not full worked solutions for these exercises.  Full integrations should therefore solve each problem from the displayed data before adding numerical conclusions here.
\end{qnbank}
